\documentclass{article}
\usepackage[margin=1in]{geometry}
\usepackage[skip=1em]{parskip}

\begin{document}
\raggedright
\textit{{\small This resume was written in LaTeX and compiled with pdflatex. See tug.org for details.}}

{\small Sample Person \\
Software Engineer \\
Email address \\
Phone number}

\underline{Education}

Sample High School (class of wxyz)
\begin{itemize}
\item Extracurriculars: Club 1, Club 2
\item Achievements: Achievement 1, Achievement 2, Achievement 3
\end{itemize}

Sample College (class of stuv)
\begin{itemize}
\item Extracurriculars: Club 1
\item Worked part-time as a X the summer before my junior year
\end{itemize}

\underline{Work experience}

Employer 1 (January 2024 - present)
\begin{itemize}
\item Brief description of work involved
\item Achievements: Achievement 1, Achievement 2, Achievement 3
\end{itemize}

Employer 2 (January 2021 - January 2023)
\begin{itemize}
\item Brief description of work involved
\item Achievements: Achievement 1, Achievement 2
\end{itemize}

\underline{Software portfolio}

Link to a software portfolio (e.g. a Github account)

\underline{Websites}

Links to any relevant websites (e.g. a chessdotcom account)

\underline{Personal statement}

This is placeholder text.

The purpose of the placeholder text is to show off the one inch margin that we set at the top of the file, with the directive that includes the geometry package, and specifies a margin of one inch all around.
\end{document}
